\section{Methodik und Vorgehensweise}
\label{sec:methodik}

\subsection{Vorgehensweise}
Die methodische Bearbeitung der Aufgabenstellung orientiert sich an einem ingenieurwissenschaftlichen Konstruktionsprozess mit den Phasen Aufgabenklärung, Konzeption, Entwurf und Ausarbeitung. Übertragen auf den Kontext Digitaler Zwillinge ergeben sich folgende Phasen:
\begin{enumerate}
    \item \textbf{Aufgabenklärung:} Analyse der Aufgabenstellung, Ableitung von Anforderungen und Erstellung einer Anforderungsliste.
    \item \textbf{Konzeption:} Recherche des Standes der Technik (Kapitel~\ref{sec:grundlagen}), Identifikation von Defiziten und Auswahl geeigneter Technologien.
    \item \textbf{Entwurf:} Modellierung der Digitalen Zwillinge in Isaac Sim, Definition der Schnittstellen zur realen Hardware.
    \item \textbf{Ausarbeitung:} Inbetriebnahme, Validierung, Reproduzierbarkeitsdokumentation und Diskussion.
\end{enumerate}

\subsection{Anforderungsanalyse}
Aus der Aufgabenstellung sowie den in Kapitel~\ref{sec:grundlagen} herausgearbeiteten Defiziten ergeben sich die in Tabelle~\ref{tab:anforderungen} dargestellten Anforderungen.

\begin{table}[htbp]
\centering
\caption{Anforderungsliste für die Anforderungsanalyse}
\label{tab:anforderungen}
\begin{tabularx}{\textwidth}{lXc}
\toprule
\textbf{Nr.} & \textbf{Anforderung} & \textbf{Klasse} \\
\midrule
A1  & Geometrische Abbildung der Anlage in OpenUSD       & Festforderung \\
A2  & Bidirektionaler Datenfluss zwischen Anlage und DT  & Festforderung \\
A3  & Physikalisch korrekte Kinematik in PhysX           & Festforderung \\
A4  & SPS-Anbindung Maze Runner via OPC UA               & Festforderung \\
A5  & Roboteranbindung Dobot via ROS\,2                  & Festforderung \\
A6  & Latenz der Datenkopplung $<$ \SI{50}{\milli\second} & Mindestforderung \\
A7  & Reproduzierbare Modellerstellung dokumentiert      & Festforderung \\
A8  & Bereitstellung als PowerPoint, Video, VR/AR        & Festforderung \\
A9  & Skalierbarkeit auf weitere Anlagen                 & Wunsch \\
A10 & Open-Access-Veröffentlichung über HsH-Bibliothek   & Wunsch \\
\bottomrule
\end{tabularx}
\end{table}

\subsection{Toolkette}
Die Bearbeitung erfolgt unter Nutzung folgender Werkzeuge:
\begin{itemize}
    \item \textbf{CAD-Aufbereitung:} SolidWorks bzw. STEP-Export der Bestandsmodelle.
    \item \textbf{Datenformat:} OpenUSD als zentrale Austauschstruktur \cite{pixar_openusd}.
    \item \textbf{Simulation:} NVIDIA Omniverse Kit mit Isaac Sim \cite{nvidia_isaacsim}.
    \item \textbf{Physik-Engine:} NVIDIA PhysX \cite{nvidia_physx}.
    \item \textbf{Middleware:} OPC UA (open62541), MQTT-Broker (Eclipse Mosquitto), ROS\,2 Humble.
    \item \textbf{Roboterbeschreibung:} URDF \cite{ros_urdf}.
    \item \textbf{Versionskontrolle:} Git mit GitLab-Instanz der HsH.
    \item \textbf{Dokumentation:} \LaTeX{} (vorliegende Arbeit), PowerPoint, OBS Studio (Videoaufzeichnung), Meta Quest 3 für VR-Demonstrationen.
\end{itemize}

\subsection{Zeit- und Arbeitsplan}
Die praktische Umsetzung beider Digitaler Zwillinge wurde ab Februar~2026 (KW~6) aufgenommen, noch vor dem offiziellen Ausgabedatum. Maze-Runner-Anlage und Dobot-Magician-Roboterarm wurden dabei parallel modelliert, erprobt und in Betrieb genommen; begleitend erfolgten Literaturrecherche und strukturelle Vorbereitung der Ausarbeitung. Die offizielle Bearbeitungszeit beträgt knapp acht Wochen ab Ausgabe am 19.06.2026~(KW~25) bis zur Abgabe am 10.08.2026~(KW~33). Tabelle~\ref{tab:zeitplan} stellt den vollständigen Projektablauf mit Arbeitspaketen und Kalenderwochen dar.

\begin{table}[htbp]
\centering
\caption{Zeit- und Arbeitsplan der Bachelorarbeit}
\label{tab:zeitplan}
\begin{tabularx}{\textwidth}{l X c}
\toprule
\textbf{AP} & \textbf{Inhalt} & \textbf{KW} \\
\midrule
\multicolumn{3}{l}{\textit{Vorbereitungsphase (Februar -- 18.\ Juni 2026)}} \\[2pt]
V1 & DZ Maze-Runner-Anlage: CAD-Aufbereitung, USD-Konvertierung, OPC-UA-Anbindung, Inbetriebnahme & 6--16 \\
V2 & DZ Dobot Magician: URDF-Modellierung, ROS\,2-Anbindung, Hardwarekonzept                     & 6--22 \\
V3 & Literaturrecherche, Theoretische Grundlagen, Strukturplanung Bachelorarbeit                  & 10--24 \\
\addlinespace
\multicolumn{3}{l}{\textit{Offizielle Bearbeitungszeit (19.\ Juni -- 10.\ August 2026, KW~25--33)}} \\[2pt]
AP1 & Aufgabenklärung, Anforderungsliste, Auswahl Toolkette                                       & 25 \\
AP2 & Ausarbeitung Kap.\,1--4 (Einleitung, Grundlagen, Methodik, Entwicklung DZ)                 & 25--28 \\
AP3 & Ausarbeitung Kap.\,5--7 (Inbetriebnahme, Potenziale, Fazit)                                & 28--31 \\
AP4 & Reproduzierbarkeitsdokumentation (PowerPoint, Video, VR/AR)                                 & 31--32 \\
AP5 & Validierung, Überarbeitung, Korrekturlesen                                                  & 32--33 \\
AP6 & Abgabe und Vorbereitung Kolloquium                                                          & 33     \\
\bottomrule
\end{tabularx}
\end{table}
